% Provenance (never printed): the money and finance strand of ../economics_series.md (grade 8: deposits and loans;
% interest as the price of borrowing; simple and compound interest; risk and duration; the interest rate after
% inflation, first look; a bank keeps part of its deposits, told). Brief: sources/book1/BRIEFS.md, g8-06.
% Chart data: figdata/grade-8/06-banks-saving-interest (JST R6 ltrate, stir, cpi; Brenland accounts).
% Ledger: g8:policy-rate-2026. The real interest rate is defined in grade 12 (present-value); here only in words.
\chapter{Banks, Saving and Interest}\label{ch:g8:banks-saving-interest}

\begin{openingchart}{Interest rate on long-term government bonds in five countries, 1870--2020, in
percent a year.}{Jord\`a--Schularick--Taylor Macrohistory Database, long-term \omterm{def:g8:banks-saving-interest:interest-rate}{interest rate}}
\begin{tikzpicture}
\begin{axis}[omtime, xmin=1870, xmax=2020, ymin=-1, ymax=21, xlabel={Year},
  ylabel={Percent a year}, extra y ticks={0}, extra tick style={grid style={black!60, thin}},
  xtick={1870,1890,1910,1930,1950,1970,1990,2010},
  legend style={at={(0.5,-0.2)}, anchor=north, legend columns=5}]
\addplot table[x=year, y=usa] {figdata/grade-8/06-banks-saving-interest/out/long_rates.csv};
\addplot table[x=year, y=gbr] {figdata/grade-8/06-banks-saving-interest/out/long_rates.csv};
\addplot table[x=year, y=fra] {figdata/grade-8/06-banks-saving-interest/out/long_rates.csv};
\addplot table[x=year, y=deu, unbounded coords=jump] {figdata/grade-8/06-banks-saving-interest/out/long_rates.csv};
\addplot table[x=year, y=ita] {figdata/grade-8/06-banks-saving-interest/out/long_rates.csv};
\legend{United States, United Kingdom, France, Germany, Italy}
\end{axis}
\end{tikzpicture}
\end{openingchart}

% figdata: facts (usa_max 13.91 1981; gbr_max 15.17 1974; fra_max 15.84 1981; deu_max 10.38 1974; ita_max 20.22
% 1982; 2020: usa 0.89, gbr 0.74, fra -0.15, deu -0.51, ita 1.17)
A government that borrows for ten years pays a yearly \omterm{def:g8:banks-saving-interest:interest-rate}{interest rate} on the \omterm{def:g7:money:money}{money}. From 1870 to 1960,
this rate was between $2.5$ and $7\,\%$ in nine years out of ten in these five countries. Then it
climbed: $13.91\,\%$ in the United States in 1981, $20.22\,\%$ in Italy in 1982. Then it fell for
forty years. In 2020, Germany and France borrowed for ten years at rates below zero: lenders accepted
to get back slightly less than they lent.

Why do people lend at all, why do rates differ, and what does a \omterm{def:g8:banks-saving-interest:bank}{bank} do with the \omterm{def:g7:money:money}{money} it receives?

\begin{recall}
\omterm{def:g7:households:saving}{Saving} is the part of \omterm{def:g7:households:income}{income} that is not spent (\cref{ch:g7:households}); \omterm{def:g7:money:money}{money} is what people
accept in payment (\cref{ch:g7:money}); the \omterm{def:g8:inflation-and-price-indices:inflation-rate}{inflation rate} is the \omterm{def:g6:reading-economic-charts:rate-of-change}{rate of change} of a \omterm{def:g8:inflation-and-price-indices:price-index}{price index}
(\cref{ch:g8:inflation-and-price-indices}); \omterm{def:g8:growth-and-compounding:compound-growth}{compound growth} multiplies a value by $(1 + g/100)$ each
year (\cref{ch:g8:growth-and-compounding}).
\end{recall}

\section{What a bank does}

\begin{definition}[Bank]\label{def:g8:banks-saving-interest:bank}
A \emph{bank}\index{bank} is a \omterm{def:g7:firms:firm}{firm} that takes \omterm{def:g8:banks-saving-interest:deposit}{deposits} and makes \omterm{def:g8:banks-saving-interest:loan}{loans}.
\end{definition}

\begin{definition}[Deposit]\label{def:g8:banks-saving-interest:deposit}
A \emph{deposit}\index{deposit} is \omterm{def:g7:money:money}{money} placed with a \omterm{def:g8:banks-saving-interest:bank}{bank}, which the \omterm{def:g8:banks-saving-interest:bank}{bank} must pay back to the
depositor, at any time or at a date agreed in advance.
\end{definition}

\begin{definition}[Loan]\label{def:g8:banks-saving-interest:loan}
A \emph{loan}\index{loan} is a sum of \omterm{def:g7:money:money}{money} lent by a lender to a borrower, who agrees to pay it back
later, usually with \omterm{def:g8:banks-saving-interest:interest}{interest}.
\end{definition}

A \omterm{def:g8:banks-saving-interest:bank}{bank} stands between people who save and people who borrow. The savers do not need to know whom to
lend to, nor check that the borrowers will repay: the \omterm{def:g8:banks-saving-interest:bank}{bank} does it for many of them at once. A family
that \omterm{def:g8:banks-saving-interest:deposit}{deposits} its \omterm{def:g7:households:saving}{savings}, a \omterm{def:g7:firms:firm}{firm} that borrows to buy a machine and a couple who borrow to buy a flat
may never meet.

\begin{omfigure}
\includegraphics[width=0.55\linewidth]{images/book1/ai/g8-banks-counter.jpg}

{\small A \omterm{def:g8:banks-saving-interest:bank}{bank} counter: a clerk hands a customer a form to open a \omterm{def:g8:banks-saving-interest:deposit}{deposit}.}
\end{omfigure}

\section{Interest, the price of borrowing}

\begin{definition}[Principal]\label{def:g8:banks-saving-interest:principal}
The \emph{principal}\index{principal} of a \omterm{def:g8:banks-saving-interest:loan}{loan} or a \omterm{def:g8:banks-saving-interest:deposit}{deposit} is the sum lent or deposited at the
start.
\end{definition}

\begin{definition}[Interest]\label{def:g8:banks-saving-interest:interest}
\emph{Interest}\index{interest} is the amount a borrower pays to a lender, on top of the \omterm{def:g8:banks-saving-interest:principal}{principal},
for the use of the \omterm{def:g7:money:money}{money}.
\end{definition}

\begin{definition}[Interest rate]\label{def:g8:banks-saving-interest:interest-rate}
The \emph{interest rate}\index{interest rate} is the \omterm{def:g8:banks-saving-interest:interest}{interest} for one year divided by the \omterm{def:g8:banks-saving-interest:principal}{principal},
in percent.
\end{definition}

The \omterm{def:g8:banks-saving-interest:interest-rate}{interest rate} is a \omterm{def:g7:money:price}{price}: the \omterm{def:g7:money:price}{price} of having \omterm{def:g7:money:money}{money} now rather than later. A borrower pays it
because a sum now is worth more to them than the same sum later: it buys a machine that earns \omterm{def:g7:money:money}{money},
or a flat to live in from today. A lender asks for it because lending means waiting, and because the
\omterm{def:g7:money:money}{money} might not come back.

\begin{measure}[A government bond's rate]
A government borrows by selling \emph{bonds}: a ten-year bond of $1\,000$ dalens promises a payment
every year and the $1\,000$ dalens back after ten years. Bonds are then bought and sold between
investors every day. When their \omterm{def:g7:money:price}{price} falls, a buyer gets the same payments for less \omterm{def:g7:money:money}{money}, so the
rate that the buyer earns rises. The ``long-term rate'' of the opening chart is that rate, worked
out from the \omterm{def:g7:markets-and-prices:market-price}{market price} of the government's ten-year bonds, one value for each year.
\end{measure}

\begin{dated}{2026}
% ledger: g8:policy-rate-2026
Central \omterm{def:g8:banks-saving-interest:bank}{banks} set the rate at which \omterm{def:g8:banks-saving-interest:bank}{banks} can place \omterm{def:g7:money:money}{money} with them or borrow from them for a short
time. Since 16 September 2026, the European Central \omterm{def:g8:banks-saving-interest:bank}{Bank} pays $2.50\,\%$ a year on the \omterm{def:g8:banks-saving-interest:deposit}{deposits} that
\omterm{def:g8:banks-saving-interest:bank}{banks} place with it overnight, up from $2.25\,\%$ in June 2026.
\end{dated}

\section{Simple and compound interest}

\begin{definition}[Simple interest]\label{def:g8:banks-saving-interest:simple-interest}
With \emph{simple interest}\index{simple interest}, the \omterm{def:g8:banks-saving-interest:interest}{interest} of each year is computed on the
\omterm{def:g8:banks-saving-interest:principal}{principal} only. After $n$ years at $r\,\%$, a \omterm{def:g8:banks-saving-interest:principal}{principal} $P$ has earned $P \times \frac{r}{100} \times
n$.
\end{definition}

\begin{definition}[Compound interest]\label{def:g8:banks-saving-interest:compound-interest}
With \emph{compound interest}\index{compound interest}, the \omterm{def:g8:banks-saving-interest:interest}{interest} of each year is added to the
\omterm{def:g8:banks-saving-interest:principal}{principal} and earns \omterm{def:g8:banks-saving-interest:interest}{interest} in the following years. After $n$ years at $r\,\%$, a \omterm{def:g8:banks-saving-interest:principal}{principal} $P$ has
become $P \times \left(1 + \frac{r}{100}\right)^{n}$.
\end{definition}

\omterm{def:g8:banks-saving-interest:compound-interest}{Compound interest} is \omterm{def:g8:growth-and-compounding:compound-growth}{compound growth} (\cref{ch:g8:growth-and-compounding}) applied to \omterm{def:g7:money:money}{money}. A \omterm{def:g8:banks-saving-interest:deposit}{deposit}
of $1\,000$ dalens at $5\,\%$:
\begin{center}
\begin{tabular}{lrrrrrr}
\toprule
After (years) & 0 & 1 & 2 & 3 & 10 & 20 \\
\midrule
\omterm{def:g8:banks-saving-interest:simple-interest}{Simple interest} & $1\,000$ & $1\,050$ & $1\,100$ & $1\,150$ & $1\,500$ & $2\,000$ \\
\omterm{def:g8:banks-saving-interest:compound-interest}{Compound interest} & $1\,000$ & $1\,050$ & $1\,102.50$ & $1\,157.63$ & $1\,628.89$ & $2\,653.30$ \\
\bottomrule
\end{tabular}
\end{center}
The gap is small in the first years and grows: after 20 years the compound \omterm{def:g8:banks-saving-interest:deposit}{deposit} has earned
$1\,653.30$ dalens of \omterm{def:g8:banks-saving-interest:interest}{interest}, against $1\,000$ with \omterm{def:g8:banks-saving-interest:simple-interest}{simple interest}. Most \omterm{def:g8:banks-saving-interest:deposit}{deposits} and \omterm{def:g8:banks-saving-interest:loan}{loans} pay
\omterm{def:g8:banks-saving-interest:compound-interest}{compound interest}.

\section{Why rates differ}

\begin{definition}[Default]\label{def:g8:banks-saving-interest:default}
A borrower \emph{defaults}\index{default} on a \omterm{def:g8:banks-saving-interest:loan}{loan} when they do not pay back the \omterm{def:g8:banks-saving-interest:interest}{interest} or the
\omterm{def:g8:banks-saving-interest:principal}{principal} on the agreed dates.
\end{definition}

\begin{model}{Risk and the interest rate}
\assumes Lenders want to be paid for waiting and for the risk that the borrower \omterm{def:g8:banks-saving-interest:default}{defaults}.
\predicts Borrowers more likely to \omterm{def:g8:banks-saving-interest:default}{default} pay higher rates; longer \omterm{def:g8:banks-saving-interest:loan}{loans} pay higher rates than short
ones.
\tested In the United States from 1870 to 2020, the ten-year government rate was above the short-term
rate in 103 of the 151 years (chart below).
\breaks When a central \omterm{def:g8:banks-saving-interest:bank}{bank} pushes short rates up or down on purpose (the United States, 1979--1981);
in panics, when short-term lenders suddenly want their \omterm{def:g7:money:money}{money} back (1873); when short-term borrowing
is in very high demand, as in the \omterm{def:g7:markets-and-prices:market}{stock-market} boom of 1928--1929; and when lenders rush
to the safest borrowers in a crisis, whose rates then fall.
\end{model}

% figdata: facts (spread_positive 103 of 151; spread_min -8.76 in 1873; 1929 -4.29; 1979 -1.75; 1980 -1.89;
% 1981 -2.48); spread (1928 -2.46). JST stir/ltrate, not printed: 1928 6.04/3.58, 1929 7.61/3.32.
\begin{omchart}{United States: the ten-year government rate minus the short-term rate, 1870--2020, in
percentage points. Above zero, lending for ten years paid more than lending for a few
months.}{Jord\`a--Schularick--Taylor Macrohistory Database, long-term and short-term \omterm{def:g8:banks-saving-interest:interest-rate}{interest rates}}
\begin{tikzpicture}
\begin{axis}[omrecord, ybar, bar width=0.6pt, xmin=1868, xmax=2022, ymin=-9.5, ymax=4,
  xtick={1870,1890,1910,1930,1950,1970,1990,2010}, xlabel={Year},
  ylabel={Percentage points}, height=0.45\linewidth,
  x tick label style={/pgf/number format/1000 sep={}}]
\addplot[fill=omDef, draw=omDef] table[x=year, y=spread] {figdata/grade-8/06-banks-saving-interest/out/spread.csv};
\end{axis}
\end{tikzpicture}
\end{omchart}

In 1873, short-term rates jumped $8.76$ points above long-term rates: in a panic, people who had lent
for a few months wanted their \omterm{def:g7:money:money}{money} back at once. In 1929 the gap was $4.29$ points, but it was
already $2.46$ points in 1928, before the \omterm{def:g7:markets-and-prices:market}{stock-market} crash of October 1929: in the boom years, much
short-term \omterm{def:g7:money:money}{money} was lent to buyers of shares, and its \omterm{def:g7:money:price}{price} stayed high. From 1979 to 1981, the central \omterm{def:g8:banks-saving-interest:bank}{bank} of the United States
raised short-term rates very high to bring inflation down.

\section{Interest after inflation, a first look}

A \omterm{def:g8:banks-saving-interest:deposit}{deposit} at $5\,\%$ a year gives $1\,050$ dalens after a year for $1\,000$ dalens deposited. If
\omterm{def:g7:money:price}{prices} rose by $5\,\%$ during the year, the $1\,050$ dalens buy exactly what $1\,000$ dalens bought a
year earlier: the saver has gained nothing. What a lender gains in goods is roughly the \omterm{def:g8:banks-saving-interest:interest-rate}{interest rate}
minus the \omterm{def:g8:inflation-and-price-indices:inflation-rate}{inflation rate}. In the United States in 1975, short-term rates were $5.82\,\%$ and inflation
$9.14\,\%$: lenders were paid less than the rise in \omterm{def:g7:money:price}{prices}.

\section{What a bank keeps at hand}

A \omterm{def:g8:banks-saving-interest:bank}{bank} does not keep all its \omterm{def:g8:banks-saving-interest:deposit}{deposits} as cash. On a normal day, only a few depositors take their \omterm{def:g7:money:money}{money}
out, so the \omterm{def:g8:banks-saving-interest:bank}{bank} keeps part of the \omterm{def:g8:banks-saving-interest:deposit}{deposits} at hand, as cash and as \omterm{def:g8:banks-saving-interest:deposit}{deposits} at the central \omterm{def:g8:banks-saving-interest:bank}{bank}, and
lends the rest.

\begin{omfigure}
\begin{tikzpicture}[font=\small]
  \node[font=\small\bfseries] at (1.5,4.4) {What the bank owes};
  \node[font=\small\bfseries] at (6.0,4.4) {What the bank has};
  \draw[fill=omDef!20] (0,0) rectangle (3,4);
  \node[align=center] at (1.5,2) {deposits\\$1\,000$};
  \draw[fill=omDef!45] (4.5,0) rectangle (7.5,3.6);
  \node[align=center] at (6.0,1.8) {loans\\$900$};
  \draw[fill=black!15] (4.5,3.6) rectangle (7.5,4);
  \node[font=\footnotesize] at (6.0,3.8) {kept at hand: $100$};
\end{tikzpicture}

{\small An invented \omterm{def:g8:banks-saving-interest:bank}{bank}: of $1\,000$ dalens of \omterm{def:g8:banks-saving-interest:deposit}{deposits}, it keeps $100$ at hand and lends $900$.}
\end{omfigure}

If too many depositors ask for their \omterm{def:g7:money:money}{money} on the same day, the \omterm{def:g8:banks-saving-interest:bank}{bank} cannot pay them all at once: its
\omterm{def:g8:banks-saving-interest:loan}{loans} come back only over years. What then happens, and how the \omterm{def:g8:banks-saving-interest:loan}{loans} of \omterm{def:g8:banks-saving-interest:bank}{banks} create new \omterm{def:g8:banks-saving-interest:deposit}{deposits},
is told in the volume for grades 10 to 12.

\section{Exercises}

\begin{exercise}[$\star$]\label{exo:g8:banks-saving-interest:1}
Use the opening chart. About how high was the United States' long-term rate at its peak, and in which
year?
\end{exercise}

\begin{exercise}[$\star$]\label{exo:g8:banks-saving-interest:2}
A \omterm{def:g8:banks-saving-interest:deposit}{deposit} of $2\,000$ dalens earns $4\,\%$ \omterm{def:g8:banks-saving-interest:simple-interest}{simple interest} a year. How much \omterm{def:g8:banks-saving-interest:interest}{interest} does it earn in
3 years?
\end{exercise}

\begin{exercise}[$\star$]\label{exo:g8:banks-saving-interest:3}
The same $2\,000$ dalens earn $4\,\%$ \omterm{def:g8:banks-saving-interest:compound-interest}{compound interest}. Compute the \omterm{def:g8:banks-saving-interest:deposit}{deposit} after 1, 2 and 3 years.
\end{exercise}

\begin{exercise}[$\star$]\label{exo:g8:banks-saving-interest:4}
Use the opening chart. In which years were the rates lowest, and which countries' rates went below
zero?
\end{exercise}

\begin{exercise}[$\star\star$]\label{exo:g8:banks-saving-interest:5}
Compute what $1\,000$ dalens become after 10 years at $5\,\%$ \omterm{def:g8:banks-saving-interest:compound-interest}{compound interest}, using a power.
\end{exercise}

\begin{exercise}[$\star\star$]\label{exo:g8:banks-saving-interest:6}
A \omterm{def:g7:firms:firm}{firm} borrows $10\,000$ dalens for one year. \omterm{def:g8:banks-saving-interest:bank}{Bank} A charges $7\,\%$ for the year. \omterm{def:g8:banks-saving-interest:bank}{Bank} B charges
$0.6\,\%$ a month, compounded each month for 12 months. Which \omterm{def:g8:banks-saving-interest:loan}{loan} costs less, and by how much?
\end{exercise}

\begin{exercise}[$\star\star$]\label{exo:g8:banks-saving-interest:7}
A \omterm{def:g8:banks-saving-interest:bank}{bank} lends to a safe \omterm{def:g7:firms:firm}{firm} at $4\,\%$. A new \omterm{def:g7:firms:firm}{firm}, which is more likely to \omterm{def:g8:banks-saving-interest:default}{default}, asks for the same
\omterm{def:g8:banks-saving-interest:loan}{loan} of $50\,000$ dalens. If the model ``Risk and the \omterm{def:g8:banks-saving-interest:interest-rate}{interest rate}'' holds, will its rate be higher or
lower? If the \omterm{def:g8:banks-saving-interest:bank}{bank} adds 2 percentage points for the risk, how much more \omterm{def:g8:banks-saving-interest:interest}{interest} does the new \omterm{def:g7:firms:firm}{firm}
pay in a year?
\end{exercise}

\begin{exercise}[$\star\star$]\label{exo:g8:banks-saving-interest:8}
In the United States, the short-term rate was $5.82\,\%$ in 1975 and $16.39\,\%$ in 1981; inflation
was $9.14\,\%$ and $9.44\,\%$. Compute the rate minus inflation in each year. In which year did
lenders gain in goods?
\end{exercise}

\begin{exercise}[$\star\star$]\label{exo:g8:banks-saving-interest:9}
Give an observation from this chapter that does not fit the model ``Risk and the \omterm{def:g8:banks-saving-interest:interest-rate}{interest rate}''.
\end{exercise}

\begin{exercise}[$\star\star\star$]\label{exo:g8:banks-saving-interest:10}
Using the rule of 70, how long does a \omterm{def:g8:banks-saving-interest:deposit}{deposit} at $5\,\%$ \omterm{def:g8:banks-saving-interest:compound-interest}{compound interest} take to double? Check with
a power.
\end{exercise}

\begin{exercise}[$\star\star\star$]\label{exo:g8:banks-saving-interest:11}
A \omterm{def:g8:banks-saving-interest:bank}{bank} keeps $10\,\%$ of its \omterm{def:g8:banks-saving-interest:deposit}{deposits} at hand. A family \omterm{def:g8:banks-saving-interest:deposit}{deposits} $4\,000$ dalens. How much of this
\omterm{def:g8:banks-saving-interest:deposit}{deposit} can the \omterm{def:g8:banks-saving-interest:bank}{bank} lend?
\end{exercise}

\begin{exercise}[$\star\star\star$]\label{exo:g8:banks-saving-interest:12}
The opening chart shows long-term rates falling from the early 1980s to 2020. Explain what the chart
alone cannot tell about why they fell.
\end{exercise}

\section{Problem: Saving for a Boat}

\begin{problem}[{Weekend problem --- a deposit at simple and compound interest, a loan at two rates,
and Brenland's interest rate after inflation}]
\label{pb:g8:banks-saving-interest:1}
Everything in this problem is invented. Mira wants a sailing boat.

\medskip
\noindent\textbf{Part I --- The deposit.} Mira \omterm{def:g8:banks-saving-interest:deposit}{deposits} $5\,000$ dalens for 10 years at $3\,\%$ a
year.
\begin{enumerate}
  \item How much \omterm{def:g8:banks-saving-interest:interest}{interest} does the \omterm{def:g8:banks-saving-interest:deposit}{deposit} earn in the first year?
  \item With \omterm{def:g8:banks-saving-interest:simple-interest}{simple interest}, how much does she have after 10 years?
  \item With \omterm{def:g8:banks-saving-interest:compound-interest}{compound interest}, how much does she have after 2 years? After 10 years?
  \item How much more does \omterm{def:g8:banks-saving-interest:compound-interest}{compound interest} give after 10 years?
\end{enumerate}

\medskip
\noindent\textbf{Part II --- The loan.} Another sailor, who has $5\,000$ dalens, wants a boat costing
$8\,000$ dalens now. He borrows the rest for 2 years, with \omterm{def:g8:banks-saving-interest:compound-interest}{compound interest}, paying everything back
at the end.
\begin{enumerate}[resume]
  \item How much does he borrow? How much does he pay back at $6\,\%$ a year?
  \item How much does he pay back at $9\,\%$ a year?
  \item How much more \omterm{def:g8:banks-saving-interest:interest}{interest} does he pay at $9\,\%$?
  \item Give a reason why one \omterm{def:g8:banks-saving-interest:bank}{bank} might offer him $6\,\%$ and another $9\,\%$.
\end{enumerate}

\medskip
\noindent\textbf{Part III --- Brenland.} Brenland's \omterm{def:g8:inflation-and-price-indices:cpi}{consumer price index} (2015 = 100) and the
central \omterm{def:g8:banks-saving-interest:bank}{bank}'s rate:
\begin{center}
\begin{tabular}{lrrrr}
\toprule
Year & 2021 & 2022 & 2024 & 2025 \\
\midrule
\omterm{def:g8:inflation-and-price-indices:cpi}{Consumer price index} & 115.9 & 127.7 & 140.0 & 143.5 \\
Central \omterm{def:g8:banks-saving-interest:bank}{bank} rate (\%) & 3.11 & 5.42 & 5.36 & 4.62 \\
\bottomrule
\end{tabular}
\end{center}
\begin{enumerate}[resume]
  \item Compute the \omterm{def:g8:inflation-and-price-indices:inflation-rate}{inflation rate} in 2022, to one decimal.
  \item Compute the central \omterm{def:g8:banks-saving-interest:bank}{bank}'s rate minus inflation in 2022. What does the sign mean for a lender?
  \item Compute the \omterm{def:g8:inflation-and-price-indices:inflation-rate}{inflation rate} and the rate minus inflation in 2025.
  \item Mira's \omterm{def:g8:banks-saving-interest:deposit}{deposit} of Part I, at \omterm{def:g8:banks-saving-interest:compound-interest}{compound interest}, is worth its 10-year value. If \omterm{def:g7:money:price}{prices} rise by
        $2.5\,\%$ a year during the 10 years, what is it worth in the \omterm{def:g7:money:price}{prices} of the start? Did her
        \omterm{def:g7:households:saving}{saving} gain in goods?
\end{enumerate}
\end{problem}
